
% Lista de abreviaturas
%--------------------------------------
\vspace*{45pt}
\begin{flushleft}
	{\Large \textbf{\scshape{List of Abbreviations}}}
\end{flushleft}
\vspace*{20pt}

\begin{description}
    \item[2D] Two-dimensional
    \item[3D] Three-dimensional
    \item[AA] Amino acid
    \item[AI] Artificial Intelligence
    \item[API] Application Programming Interface
    \item[CNN] Convolutional Neural Network
    \item[CPU] Central Processing Unit
    \item[cryo-EM] Cryogenic Electron Microscopy
    \item[CSS] Cascading Style Sheets
    \item[CUDA] Compute Unified Device Architecture
    \item[DB] Database
    \item[DQN] Deep Q-Network
    \item[DRL] Deep Reinforcement Learning
    \item[FASTA] FAST-All sequence format
    \item[GO] Gene Ontology
    \item[GO:BP] Gene Ontology: Biological Process
    \item[GO:CC] Gene Ontology: Cellular Component
    \item[GO:MF] Gene Ontology: Molecular Function
    \item[GET] HTTP GET request method
    \item[GPU] Graphics Processing Unit
    \item[H] Hydrophobic residue
    \item[HC] Hill Climbing
    \item[HP] Hydrophobic-Polar
    \item[HTML] HyperText Markup Language
    \item[IEEE] \textit{Institute of Electrical and Electronics Engineers}
    \item[ISEC] Instituto Superior de Engenharia de Coimbra
    \item[JSON] JavaScript Object Notation
    \item[MC] Monte Carlo
    \item[MDP] Markov Decision Process
    \item[MSE] Mean Squared Error
    \item[MVT] Model-View-Template
    \item[NP] Nondeterministic Polynomial time
    \item[P] Polar residue
    \item[PDB] Protein Data Bank
    \item[POST] HTTP POST request method
    \item[PPO] Proximal Policy Optimization
    \item[QL] Q-Learning
    \item[REMC] Replica Exchange Monte Carlo
    \item[RL] Reinforcement Learning
    \item[RMSD] Root Mean Square Deviation
    \item[SA] Simulated Annealing
    \item[SAC] Soft Actor-Critic
    \item[SAW] Self-Avoiding Walk
    \item[SVG] Scalable Vector Graphics
    \item[TD] Temporal Difference
    \item[TM-score] Template Modelling score
    \item[UniProt] Universal Protein Resource
\end{description}
% !TEX root = ../main.tex
